\section{2024 Probability and Statistics}

\begin{exercise}\label{probability:2024:1}
The sequence $(X_1, X_2, \ldots, X_n, \ldots)$ is a Dirichlet process with base distribution $G_0$ and concentration parameter $\alpha_0 > 0$.
(a) Derive the distribution of $X_n, n \geq 1$.
(b) Let $Y_n = I(X_n > 0)$. Prove or disprove $(Y_n)_{n \geq 1}$ forms a Dirichlet Process.
\end{exercise}

\begin{solution}
\textbf{Prerequisites / Core Concepts:}
\begin{itemize}
    \item \textbf{Dirichlet Process (DP) / Pólya Urn Scheme:} A sequence of random variables $X_1, X_2, \dots$ drawn from a Dirichlet Process $DP(\alpha_0, G_0)$ can be described sequentially. The first draw is from the base distribution $G_0$. Subsequent draws are either identical to a previously drawn value (with probability proportional to how many times it was drawn) or a fresh draw from $G_0$.
    \item \textbf{Conditional Probability and Marginals:} The marginal distribution of a variable can be found by taking the expectation of its conditional distribution: $\mathbb{P}(X \in A) = \mathbb{E}[\mathbb{P}(X \in A \mid Y)]$.
\end{itemize}

\textbf{Intuition / Motivation:}

(a) Imagine a magic urn. You draw a ball. With some probability, you paint it a brand new color from a catalog ($G_0$). Otherwise, you copy the color of a ball already in the urn.
If you look at the $n$-th ball drawn, what is the probability it is a specific color? Even though the *conditional* probability heavily depends on the colors of the first $n-1$ balls, the *unconditional* (marginal) probability before the game starts is perfectly symmetric. Every ball, whether it's the 1st or the 100th, has the exact same marginal distribution $G_0$. The rich-get-richer dynamics only affect the joint distribution, not the marginals.
(b) Now, instead of looking at the exact color, we just ask: "Is the number drawn greater than 0?" This turns our complex continuous process into a simple sequence of Yes/No (1/0) answers.
Does this simplified Yes/No sequence still behave like a Dirichlet Process? Yes! The "rich-get-richer" copying mechanism of the Dirichlet Process is perfectly preserved when you lump categories together. If you group all positive numbers into "Yes" and all negative numbers into "No", the rules of drawing a new "Yes" or copying an old "Yes" match exactly with a new Dirichlet Process whose base distribution is just a coin flip.

\textbf{Logical Step-by-Step Breakdown:}

\textbf{Part (a): Marginal distribution of $X_n$}

\textbf{Step 1: Base case for induction.}
By the definition of the Pólya Urn representation of a Dirichlet Process, the first variable is drawn directly from the base distribution:
\[ X_1 \sim G_0 \]
Thus, $\mathbb{P}(X_1 \in A) = G_0(A)$ for any measurable set $A$.

\textbf{Step 2: Set up the inductive step.}
Assume that for all $i \leq n$, the marginal distribution is $G_0$, meaning $\mathbb{P}(X_i \in A) = G_0(A)$.
We want to find the marginal distribution of $X_{n+1}$. The predictive distribution of $X_{n+1}$ given the past sequence is:
\[ \mathbb{P}(X_{n+1} \in A \mid X_1, \dots, X_n) = \frac{\alpha_0}{\alpha_0 + n} G_0(A) + \frac{1}{\alpha_0 + n} \sum_{i=1}^n \delta_{X_i}(A) \]
where $\delta_{X_i}(A)$ is the indicator function $\mathbb{I}(X_i \in A)$.

\textbf{Step 3: Calculate the marginal probability.}
We take the expectation of the conditional probability over the joint distribution of $(X_1, \dots, X_n)$:
\[ \mathbb{P}(X_{n+1} \in A) = \mathbb{E} \left[ \mathbb{P}(X_{n+1} \in A \mid X_1, \dots, X_n) \right] \]
\[ = \frac{\alpha_0}{\alpha_0 + n} G_0(A) + \frac{1}{\alpha_0 + n} \sum_{i=1}^n \mathbb{E}[\delta_{X_i}(A)] \]
Since $\mathbb{E}[\delta_{X_i}(A)] = \mathbb{P}(X_i \in A)$, and by our inductive hypothesis this equals $G_0(A)$, we substitute:
\[ \mathbb{P}(X_{n+1} \in A) = \frac{\alpha_0}{\alpha_0 + n} G_0(A) + \frac{1}{\alpha_0 + n} \sum_{i=1}^n G_0(A) \]
\[ = \frac{\alpha_0 G_0(A) + n G_0(A)}{\alpha_0 + n} = \frac{(\alpha_0 + n) G_0(A)}{\alpha_0 + n} = G_0(A) \]
Thus, by mathematical induction, $X_n \sim G_0$ for all $n \geq 1$.

\textbf{Part (b): Dirichlet Process property of $Y_n$}

\textbf{Step 4: Define the new base distribution.}
Let $A = (0, \infty)$. We are defining $Y_n = \mathbb{I}(X_n > 0) = \mathbb{I}(X_n \in A)$.
Since $X_n \sim G_0$, the marginal distribution of $Y_1$ is a Bernoulli variable:
\[ \mathbb{P}(Y_1 = 1) = \mathbb{P}(X_1 \in A) = G_0(A) = p_0 \]
Thus, the new base distribution is $H_0 = \text{Bernoulli}(p_0)$.

\textbf{Step 5: Verify the predictive distribution matches a DP.}
To prove $\{Y_n\}$ forms a Dirichlet Process $DP(\alpha_0, H_0)$, we must show its predictive distribution follows the standard Pólya urn formula for $H_0$.
Using the predictive distribution of $X_{n+1}$:
\[ \mathbb{P}(Y_{n+1} = 1 \mid X_1, \dots, X_n) = \mathbb{P}(X_{n+1} \in A \mid X_1, \dots, X_n) \]
\[ = \frac{\alpha_0 G_0(A) + \sum_{i=1}^n \delta_{X_i}(A)}{\alpha_0 + n} \]
Notice that $G_0(A) = p_0$ and $\delta_{X_i}(A) = \mathbb{I}(X_i > 0) = Y_i$. Substituting these in:
\[ \mathbb{P}(Y_{n+1} = 1 \mid X_1, \dots, X_n) = \frac{\alpha_0 p_0 + \sum_{i=1}^n Y_i}{\alpha_0 + n} \]
Because this conditional probability depends on the history $(X_1, \dots, X_n)$ *only* through the sequence $(Y_1, \dots, Y_n)$, we can safely write:
\[ \mathbb{P}(Y_{n+1} = 1 \mid Y_1, \dots, Y_n) = \frac{\alpha_0 p_0 + \sum_{i=1}^n Y_i}{\alpha_0 + n} \]
This is exactly the Pólya urn predictive update rule for a Dirichlet Process with concentration parameter $\alpha_0$ and base distribution $\text{Bernoulli}(p_0)$.
Therefore, we prove that $(Y_n)_{n \geq 1}$ does indeed form a Dirichlet Process.
\end{solution}

\begin{exercise}\label{probability:2024:2}
$U_n = \sum_{j=1}^{2^n} \frac{j-1}{2^n} I\left(\frac{j-1}{2^n} \leq \phi(X) < \frac{j}{2^n}\right)$, $V_n = E(H(X,Y) | U_n)$.
Prove AS convergence of $V_n$.
\end{exercise}

\begin{solution}
\textbf{Prerequisites / Core Concepts:}
\begin{itemize}
    \item \textbf{Filtrations and $\sigma$-algebras:} A $\sigma$-algebra $\mathcal{G}_n$ represents the "information available" at time $n$. An increasing sequence of $\sigma$-algebras $\mathcal{G}_1 \subset \mathcal{G}_2 \subset \dots$ is called a filtration.
    \item \textbf{Lévy's Upward Theorem (Martingale Convergence):} If you have a random variable $W$ with finite expectation ($\mathbb{E}|W| < \infty$), the sequence of conditional expectations $V_n = \mathbb{E}[W \mid \mathcal{G}_n]$ forms a uniformly integrable martingale. It converges almost surely to $\mathbb{E}[W \mid \mathcal{G}_\infty]$, where $\mathcal{G}_\infty$ is the $\sigma$-algebra generated by the union of all $\mathcal{G}_n$.
\end{itemize}

\textbf{Intuition / Motivation:}

Imagine trying to guess the value of a function $H(X,Y)$, but you only get a blurred, pixelated image of $X$.
$U_n$ represents this pixelated version of $X$. We chop the interval $[0,1)$ into $2^n$ tiny bins. $U_n$ tells us exactly which bin $\phi(X)$ fell into.
As $n$ increases, we double the number of bins. The pixels get smaller and smaller, so the image gets sharper and sharper. The information we have at step $n$ (represented by $\mathcal{G}_n$) completely contains the information we had at step $n-1$.
$V_n$ is our "best guess" for $H(X,Y)$ given the pixelated image $U_n$. Because our sequence of guesses is based on strictly improving information, the sequence $V_n$ forms a martingale.
By Lévy's Upward Theorem, as the pixels become infinitely small ($n \to \infty$), our sequence of guesses must converge to the best guess we could make if we had a perfectly sharp image of $X$. Because the infinitely fine bins perfectly pinpoint $\phi(X)$, and $\phi$ is an invertible function, knowing the infinitely fine bins is exactly equivalent to knowing the exact value of $X$.

\textbf{Logical Step-by-Step Breakdown:}

\textbf{Step 1: Identify the information structure ($\sigma$-fields).}

Let $\mathcal{G}_n = \sigma(U_n)$ be the $\sigma$-algebra generated by the random variable $U_n$.
The variable $U_n$ discretizes $\phi(X)$ into dyadic intervals of length $1/2^n$: $[0, 1/2^n), [1/2^n, 2/2^n), \dots$
At step $n+1$, the intervals are of length $1/2^{n+1}$. Notice that every interval at step $n$ is perfectly split into two smaller intervals at step $n+1$.
Because $U_{n+1}$ provides a strictly finer partition than $U_n$, knowing $U_{n+1}$ means you automatically know $U_n$.
Therefore, $\mathcal{G}_n \subseteq \mathcal{G}_{n+1}$. The sequence $\{\mathcal{G}_n\}_{n \geq 1}$ forms a valid filtration.

\textbf{Step 2: Apply Lévy's Upward Theorem.}

We are given the sequence of random variables $V_n = \mathbb{E}[H(X,Y) \mid U_n] = \mathbb{E}[H(X,Y) \mid \mathcal{G}_n]$.
By the tower property of conditional expectation, $V_n$ is a Doob martingale with respect to the filtration $\{\mathcal{G}_n\}$.
Assuming $\mathbb{E}[|H(X,Y)|] < \infty$ (a standard implicit requirement for conditional expectations to be well-behaved), this martingale is uniformly integrable.
By Lévy's Upward Theorem (a consequence of the Martingale Convergence Theorem), $V_n$ converges almost surely to the conditional expectation given the limit of the $\sigma$-algebras:
\[ V_n \xrightarrow{\text{a.s.}} \mathbb{E}[H(X,Y) \mid \mathcal{G}_\infty] \]
where $\mathcal{G}_\infty = \sigma\left(\bigcup_{n=1}^\infty \mathcal{G}_n \right)$.

\textbf{Step 3: Characterize the limiting $\sigma$-algebra $\mathcal{G}_\infty$.}

What information is contained in $\mathcal{G}_\infty$? It contains the knowledge of which dyadic interval $\phi(X)$ falls into for *any* arbitrarily small resolution.
The dyadic intervals generate the standard Borel $\sigma$-algebra on $[0, 1)$. Therefore, knowing which interval $\phi(X)$ belongs to for all $n$ uniquely identifies the exact real value of $\phi(X)$. Thus, $\mathcal{G}_\infty = \sigma(\phi(X))$.
Furthermore, the problem context implies $\phi(u)$ (e.g., $\phi(u) = u/(1+u)$ mapping $[0, \infty) \to [0, 1)$) is a bijection. Because $\phi$ is strictly monotonic and invertible, knowing $\phi(X)$ is mathematically identical to knowing $X$.
Therefore, $\sigma(\phi(X)) = \sigma(X)$.
We conclude that $\mathcal{G}_\infty = \sigma(X)$.

\textbf{Step 4: Final Conclusion.}

Substituting $\sigma(X)$ into the limit from Step 2, we get:
\[ V_n \xrightarrow{\text{a.s.}} \mathbb{E}[H(X,Y) \mid \sigma(X)] \]
Which is standardly written as:
\[ V_n \xrightarrow{\text{a.s.}} \mathbb{E}[H(X,Y) \mid X] \]
This proves the almost sure convergence of the sequence $V_n$.
\end{solution}

\begin{exercise}\label{probability:2024:3}
Let $X$ have a uniform distribution on $[0,1]$ and let $N_{m,k}$ be the digit in the $m$th place to the right of the decimal point in $X^k$.
\begin{enumerate}
    \item[(a)] Find $\lim_{m\to\infty}\mathbb{P}(N_{m,m}=i)$ for $i=0,1,\ldots,9$.
    \item[(b)] Let $k(m)>1$. Find a necessary and sufficient condition on $k(m)$ such that
    \[
    \lim_{m\to\infty}\mathbb{P}(N_{m,k(m)}=i)=\frac{1}{10},
    \qquad i=0,1,\ldots,9.
    \]
\end{enumerate}
\end{exercise}

\begin{solution}
Put $N=10^m$ and $a=1/k$. Since $X^k$ has distribution function $y^a$ on $[0,1]$,
\[
\mathbb{P}(N_{m,k}=i)
=N^{-a}\sum_{\substack{0\leq q\leq N-1\\ q\equiv i\pmod {10}}}
\left((q+1)^a-q^a\right),
\]
where $0^a=0$. This formula is just the partition of $[0,1]$ into the intervals on which $\lfloor 10^mX^k\rfloor\equiv i\pmod {10}$.

For part (a), take $k=m$, so $a=1/m$ and $N^a=10$. The interval $0\leq X^m<10^{-m}$ contributes
\[
\mathbb{P}(X^m<10^{-m})=\mathbb{P}(X<1/10)=\frac{1}{10}
\]
to the digit $0$ and to no other digit. On the remaining intervals, the increments of $x^{1/m}$ over the ten residue classes are asymptotically equal. More explicitly, for
\[
d_q=(q+1)^{1/m}-q^{1/m},
\]
the difference between two residue-class sums over each complete block of length $10$ is bounded by the variation of the decreasing sequence $d_q$ over that block; summing over all blocks gives an $o(1)$ difference as $m\to\infty$. Since
\[
\sum_{q=1}^{10^m-1}d_q=10-1=9,
\]
each residue class receives limiting increment $9/10$ from $q\geq1$. Multiplying by $N^{-1/m}=1/10$, we obtain
\[
\lim_{m\to\infty}\mathbb{P}(N_{m,m}=0)=\frac{1}{10}+\frac{9}{100}=\frac{19}{100},
\]
and, for $i=1,\ldots,9$,
\[
\lim_{m\to\infty}\mathbb{P}(N_{m,m}=i)=\frac{9}{100}.
\]

For part (b), let $a_m=1/k(m)$ and $L_m=10^{m/k(m)}=N^{a_m}$. The atom-like initial contribution to the digit $0$ is
\[
\mathbb{P}(X^{k(m)}<10^{-m})=N^{-a_m}=L_m^{-1}.
\]
The same block-variation argument shows that the remaining mass is asymptotically split equally among the ten residue classes whenever $L_m\to\infty$. Hence all ten limiting digit probabilities are $1/10$ if
\[
L_m\to\infty,
\]
which is equivalent to
\[
\frac{m}{k(m)}\to\infty.
\]
This condition is also necessary. If $m/k(m)$ does not tend to infinity, then along a subsequence $L_m$ stays bounded, and the digit $0$ keeps the additional contribution $L_m^{-1}$ from the interval $[0,10^{-m})$; the limiting distribution cannot be uniform over the ten digits.
\end{solution}

\begin{exercise}\label{probability:2024:4}
Prove $P(A|T) = Q(T,A)$ almost surely.
\end{exercise}

\begin{solution}
\textbf{Prerequisites / Core Concepts:}
\begin{itemize}
    \item \textbf{Conditional Expectation as a Random Variable:} In measure-theoretic probability, the conditional probability $\mathbb{P}(A \mid T)$ is defined as the conditional expectation $\mathbb{E}[\mathbb{I}_A \mid T]$. It is a random variable that must satisfy two strict conditions: it must be measurable with respect to the information provided by $T$, and it must preserve probabilities when averaged over any set defined by $T$.
    \item \textbf{Regular Conditional Probability:} A function $Q(t, A)$ that acts like a valid probability measure for a fixed $t$, and for a fixed event $A$, $Q(T(\omega), A)$ acts as the conditional probability $\mathbb{P}(A \mid T)$.
    \item \textbf{Change of Variables / Pushforward Measure:} Integrating a function of a random variable $T$ with respect to the underlying probability measure $\mathbb{P}$ is equivalent to integrating the function with respect to the distribution of $T$ (often denoted $\Pi$ or $P_T$).
\end{itemize}

\textbf{Intuition / Motivation:}

How do we prove that a specific formula $Q(T,A)$ is actually the true conditional probability $\mathbb{P}(A \mid T)$?
In advanced probability, conditional probability isn't just a simple fraction; it's a rule that assigns a probability to event $A$ based on the exact observed value of $T$.
To prove $Q(T,A)$ is the correct rule, we must put it through the "Partial Averaging Test" (the definition of conditional expectation).
If $Q(T,A)$ is the true rule, then if we look at any specific subset of outcomes for $T$ (let's say, $T \in [1, 5]$) and average our rule $Q(T,A)$ over all those specific outcomes, we must recover the exact unconditional probability of the joint event "$A$ occurs AND $T \in [1, 5]$".
By using the change of variables formula, we can seamlessly translate the integral over the abstract sample space into a standard integral over the real numbers using the distribution of $T$, confirming that the rule perfectly satisfies the rigorous definition.

\textbf{Logical Step-by-Step Breakdown:}

\textbf{Step 1: State the definition of Conditional Probability.}
By the Radon-Nikodym theorem, a random variable $W$ is the conditional probability $\mathbb{P}(A \mid T)$ (which is $\mathbb{E}[\mathbb{I}_A \mid \sigma(T)]$) almost surely if it satisfies two conditions:
1. $W$ is $\sigma(T)$-measurable.
2. For any event $B \in \sigma(T)$, the partial expectation matches: $\int_B W d\mathbb{P} = \mathbb{P}(A \cap B)$.

\textbf{Step 2: Verify Measurability.}
Let $W = Q(T, A)$.
By definition, a regular conditional probability $Q(t, A)$ is a Borel-measurable function of $t$ for any fixed $A$.
Since $T$ is a random variable (measurable mapping from the sample space $\Omega$ to $\mathbb{R}$) and $Q(\cdot, A)$ is Borel-measurable, their composition $W(\omega) = Q(T(\omega), A)$ is $\sigma(T)$-measurable.
This satisfies the first condition.

\textbf{Step 3: Set up the Partial Averaging Test.}
Let $B \in \sigma(T)$. By the definition of $\sigma(T)$, there exists a Borel set $D \in \mathcal{B}(\mathbb{R})$ such that $B = \{ \omega : T(\omega) \in D \}$.
We must evaluate the integral of $W$ over the set $B$:
\[ \int_B Q(T, A) d\mathbb{P} = \int_\Omega Q(T(\omega), A) \mathbb{I}_D(T(\omega)) \mathbb{P}(d\omega) \]

\textbf{Step 4: Apply the Change of Variables Formula.}
We use the Law of the Unconscious Statistician (or change of variables for pushforward measures). Let $\Pi = \mathbb{P} \circ T^{-1}$ be the probability distribution of $T$.
We can rewrite the integral over $\Omega$ as an integral over the real line $\mathbb{R}$:
\[ \int_\Omega Q(T(\omega), A) \mathbb{I}_D(T(\omega)) \mathbb{P}(d\omega) = \int_{\mathbb{R}} Q(t, A) \mathbb{I}_D(t) \Pi(dt) = \int_D Q(t, A) \Pi(dt) \]
The problem statement implies that $Q(t, A)$ defines the conditional probability such that averaging $Q$ over a set $D$ yields the joint probability: $\int_D Q(t,A) \Pi(dt) = \mathbb{P}(A \cap \{T \in D\})$.
Since $B = \{T \in D\}$, this equals $\mathbb{P}(A \cap B)$.

\textbf{Step 5: Conclude.}
We have shown that $\int_B Q(T, A) d\mathbb{P} = \mathbb{P}(A \cap B)$ for all $B \in \sigma(T)$.
Because $Q(T, A)$ satisfies both properties of conditional expectation, it must be the conditional probability. Since conditional expectations are unique up to sets of measure zero, we conclude:
\[ \mathbb{P}(A \mid T) = Q(T, A) \quad \text{almost surely.} \]
\end{solution}

\begin{exercise}\label{probability:2024:5}
Consider a random sample of size $n$, written as an $r=r_n$ by $c=c_n$ matrix $\{X_{ij}:1\leq i\leq r_n,\ 1\leq j\leq c_n\}$ with $n=r_nc_n$. The $X_{ij}$ are i.i.d. with c.d.f. $F$ and continuous density $f$. Let $\beta$ be the median, so $F(\beta)=1/2$, and define
\[
\widehat{\beta}_n=\min_{1\leq j\leq c_n}\max_{1\leq i\leq r_n}X_{ij}.
\]
\begin{enumerate}
    \item[(a)] What condition on $r_n$ gives median-unbiasedness asymptotically, i.e. $\beta$ is the median of the limiting distribution of $\widehat{\beta}_n$?
    \item[(b)] If $F$ is differentiable near $\beta$ and $f(\beta)>0$, find the limiting distribution of $r_n(\widehat{\beta}_n-\beta)$ under the condition in part (a).
\end{enumerate}
\end{exercise}

\begin{solution}
For one column, let
\[
M_j=\max_{1\leq i\leq r}X_{ij}.
\]
Then
\[
\mathbb{P}(M_j\leq x)=F(x)^r.
\]
The column maxima are independent across $j$, and $\widehat{\beta}_n=\min_j M_j$. Hence
\[
\mathbb{P}(\widehat{\beta}_n\leq x)
=1-\mathbb{P}(M_1>x,\ldots,M_c>x)
=1-\left(1-F(x)^r\right)^c.
\]
At $x=\beta$ this gives
\[
\mathbb{P}(\widehat{\beta}_n\leq \beta)
=1-\left(1-2^{-r}\right)^c.
\]
Therefore the asymptotic median-unbiasedness condition is
\[
\left(1-2^{-r_n}\right)^{c_n}\to\frac{1}{2}.
\]
Since $c_n=n/r_n$ and $\log(1-u)\sim -u$ as $u\downarrow0$, this is equivalently
\[
\frac{n}{r_n}2^{-r_n}\to \log 2.
\]

Now set
\[
x_n(t)=\beta+\frac{t}{r_n}.
\]
Using differentiability at $\beta$,
\[
F(x_n(t))
=\frac{1}{2}+\frac{f(\beta)t}{r_n}+o(r_n^{-1})
=\frac{1}{2}\left(1+\frac{2f(\beta)t}{r_n}+o(r_n^{-1})\right).
\]
Thus
\[
F(x_n(t))^{r_n}
=2^{-r_n}\left(1+\frac{2f(\beta)t}{r_n}+o(r_n^{-1})\right)^{r_n}
\to 2^{-r_n}e^{2f(\beta)t}
\]
in the sense needed after multiplication by $c_n$. By the condition above,
\[
c_nF(x_n(t))^{r_n}\to (\log 2)e^{2f(\beta)t}.
\]
Therefore
\[
\begin{aligned}
\mathbb{P}\{r_n(\widehat{\beta}_n-\beta)\leq t\}
&=1-\left(1-F(x_n(t))^{r_n}\right)^{c_n}\\
&\to 1-\exp\left\{-(\log 2)e^{2f(\beta)t}\right\}.
\end{aligned}
\]
The limiting distribution function is
\[
\boxed{
G(t)=1-\exp\left\{-(\log 2)e^{2f(\beta)t}\right\}
=1-2^{-e^{2f(\beta)t}}.
}
\]
At $t=0$, $G(0)=1/2$, matching the median normalization.
\end{solution}

\begin{exercise}\label{probability:2024:6}
Poisson hidden parameter. $\hat{\theta}_n = \bar{Y}_n / (2\bar{X}_n)$.
(a) Asymptotic distribution.
(b) Exact test for $\theta=1/2$.
\end{exercise}

\begin{solution}
\textbf{Prerequisites / Core Concepts:}
\begin{itemize}
    \item \textbf{Delta Method:} A statistical technique to derive the asymptotic normal distribution of a function of random variables. If $\sqrt{n}(\bar{V} - \mu) \xrightarrow{d} N(0, \Sigma)$, then $\sqrt{n}(g(\bar{V}) - g(\mu)) \xrightarrow{d} N(0, \nabla g(\mu)^T \Sigma \nabla g(\mu))$.
    \item \textbf{Law of Total Expectation / Variance:} $\mathbb{E}[X] = \mathbb{E}[\mathbb{E}[X \mid S]]$. Useful for dealing with hierarchical models.
    \item \textbf{Exact Tests:} A statistical test where the p-value can be calculated exactly using the true discrete distribution (like a Binomial), rather than relying on large-sample continuous approximations (like a Normal distribution).
\end{itemize}

\textbf{Intuition / Motivation:}

We have two variables, $X$ and $Y$, generated from a complex hidden process involving a Poisson variable $S$. We are using their sample averages to estimate a parameter $\theta$ via the ratio $\hat{\theta}_n = \bar{Y}_n / (2\bar{X}_n)$.
(a) To find how this estimate behaves when we have a huge amount of data ($n \to \infty$), we don't need to unravel the whole complex Poisson distribution. We just need the expected values and variances of $X$ and $Y$. By calculating these moments using the Law of Total Expectation, we can use the Delta Method. The Delta Method acts like a Taylor Expansion for probability: it tells us that a ratio of normal-looking averages is itself normally distributed, and provides a formula for its variance based on the gradients.
(b) If we want to test if $\theta$ is exactly $1/2$, the math suddenly collapses into something beautiful. If $\theta=1/2$, $X$ and $Y$ become perfectly symmetric clones of each other! If they are clones, then telling them apart is exactly like flipping a fair coin. Given the total number of "successes" across both $X$ and $Y$, the number belonging to $Y$ should simply follow a Binomial distribution with $p=0.5$. This allows us to run a perfect, exact test without needing large $n$.

\textbf{Logical Step-by-Step Breakdown:}

\textbf{Part (a): Asymptotic distribution using the Delta Method}

\textbf{Step 1: Calculate the first moments of $X$ and $Y$.}
We use the Law of Total Expectation conditioning on $S \sim \text{Poisson}(2\lambda)$. The variables $X$ and $Y$ are given conditionally as indicators with $P(X=1|S=s) = 1/2^{s+1}$ and $P(Y=1|S=s) = \theta/2^s$.
\[ \mathbb{E}[X] = \mathbb{E}[\mathbb{E}[X \mid S]] = \sum_{s=0}^\infty \frac{1}{2^{s+1}} \frac{e^{-2\lambda} (2\lambda)^s}{s!} = \frac{e^{-2\lambda}}{2} \sum_{s=0}^\infty \frac{\lambda^s}{s!} = \frac{1}{2} e^{-2\lambda} e^\lambda = \frac{1}{2} e^{-\lambda} \]
\[ \mathbb{E}[Y] = \mathbb{E}[\mathbb{E}[Y \mid S]] = \sum_{s=0}^\infty \frac{\theta}{2^s} \frac{e^{-2\lambda} (2\lambda)^s}{s!} = \theta e^{-2\lambda} \sum_{s=0}^\infty \frac{\lambda^s}{s!} = \theta e^{-\lambda} \]

\textbf{Step 2: Calculate the second moments and covariance.}
Since $X$ and $Y$ are binary indicator variables (implied by the conditional probabilities), $X^2 = X$ and $Y^2 = Y$.
Thus, $\text{Var}(X) = \mathbb{E}[X] - \mathbb{E}[X]^2 = \frac{1}{2} e^{-\lambda} - \frac{1}{4} e^{-2\lambda}$.
And $\text{Var}(Y) = \mathbb{E}[Y] - \mathbb{E}[Y]^2 = \theta e^{-\lambda} - \theta^2 e^{-2\lambda}$.
For the covariance, assuming $X$ and $Y$ are conditionally independent given $S$:
\[ \mathbb{E}[XY] = \mathbb{E}[\mathbb{E}[X \mid S]\mathbb{E}[Y \mid S]] = \mathbb{E}\left[ \frac{1}{2^{S+1}} \frac{\theta}{2^S} \right] = \frac{\theta}{2} \mathbb{E}\left[ 4^{-S} \right] \]
\[ \mathbb{E}[XY] = \frac{\theta}{2} \sum_{s=0}^\infty 4^{-s} \frac{e^{-2\lambda} (2\lambda)^s}{s!} = \frac{\theta}{2} e^{-2\lambda} \sum_{s=0}^\infty \frac{(\lambda/2)^s}{s!} = \frac{\theta}{2} e^{-2\lambda} e^{\lambda/2} = \frac{\theta}{2} e^{-1.5\lambda} \]
$\text{Cov}(X,Y) = \mathbb{E}[XY] - \mathbb{E}[X]\mathbb{E}[Y] = \frac{\theta}{2} e^{-1.5\lambda} - \frac{\theta}{2} e^{-2\lambda}$.

\textbf{Step 3: Apply the Delta Method.}
Let the function be $g(x, y) = \frac{y}{2x}$. We are analyzing $g(\bar{X}_n, \bar{Y}_n)$.
By the Central Limit Theorem, the sample means are asymptotically normal.
We evaluate the gradient of $g$ at the true means $\mu_x = \frac{1}{2} e^{-\lambda}$ and $\mu_y = \theta e^{-\lambda}$:
\[ \frac{\partial g}{\partial x} = -\frac{y}{2x^2} \implies \text{at means: } -\frac{\theta e^{-\lambda}}{2(1/4 e^{-2\lambda})} = -2\theta e^\lambda \]
\[ \frac{\partial g}{\partial y} = \frac{1}{2x} \implies \text{at means: } \frac{1}{2(1/2 e^{-\lambda})} = e^\lambda \]
The asymptotic variance is $V = \nabla g^T \Sigma \nabla g$, where $\Sigma$ is the covariance matrix.
\[ V = \left(-2\theta e^\lambda\right)^2 \text{Var}(X) + \left(e^\lambda\right)^2 \text{Var}(Y) + 2(-2\theta e^\lambda)(e^\lambda)\text{Cov}(X,Y) \]
By the Delta method, the asymptotic distribution is:
\[ \sqrt{n}(\hat{\theta}_n - \theta) \xrightarrow{d} \mathcal{N}(0, V) \]

\textbf{Part (b): Exact test for $\theta = 1/2$}

\textbf{Step 4: Establish symmetry under the null hypothesis.}
Under the null hypothesis $H_0: \theta = 1/2$, substitute $\theta$ into the conditional probabilities:
\[ \mathbb{P}(X=1 \mid S=s) = \frac{1}{2^{s+1}} \]
\[ \mathbb{P}(Y=1 \mid S=s) = \frac{1/2}{2^s} = \frac{1}{2^{s+1}} \]
Because their conditional probabilities are identical, $X$ and $Y$ are fully symmetric and unconditionally identically distributed Bernoulli random variables with $p = \frac{1}{2} e^{-\lambda}$.

\textbf{Step 5: Formulate the exact Binomial test.}
Let $N_X = \sum_{i=1}^n X_i$ and $N_Y = \sum_{i=1}^n Y_i$. Under $H_0$, $N_X$ and $N_Y$ are drawn from the exact same Poisson-Binomial mixture distribution.
To remove the dependence on the nuisance parameter $\lambda$, we condition on the total number of successes $K = N_X + N_Y$.
Given a fixed total of $K$ independent, identically distributed successes, the assignment of any individual success to sequence $Y$ rather than $X$ is essentially a fair coin flip.
Therefore, conditioned on $N_X + N_Y = K$, the distribution of $N_Y$ is exactly Binomial:
\[ N_Y \mid (N_X + N_Y = K) \sim \text{Binomial}\left(K, \frac{1}{2}\right) \]
An exact test for $H_0: \theta = 1/2$ can be performed by computing the $p$-value directly from this $\text{Binomial}(K, 0.5)$ distribution, evaluating the probability of observing a value as extreme as the actual $N_Y$.
\end{solution}
